\documentclass[12pt]{article}
\usepackage[utf8]{inputenc}
\usepackage[T1]{fontenc}
\usepackage{hyperref}
\title{Yameo Language Resources for K-12 Curriculum Development}
\author{Compiled by Claude Sonnet 5.5 (claude-sonnet-5-5) for Andr\'as Kornai}
\date{}
\begin{document}
\maketitle

\section{Introduction}
Yameo (ISO 639-3: yme; Glottocode yame1242) is an extinct language of the Peba-Yaguan family. It was spoken along the Amazon in Loreto, Peru, between the Tigre and Nanay rivers. Wikipedia versions disagree on when it died out, giving either after 1964 or the late 20th century. Wikipedia cites a 1950 account by John Alden Mason, which includes a list of dialect groups. Yameo's only closely related documented languages are Peba, which is also extinct, and Yagua, which is still spoken. Masamae, spoken near the Mazan River, is closely related to Yameo or a variety of it, depending on the source.

Nearly all of the search results were Wikipedia mirrors. The only primary-type linguistic data found was a numeral list (1 to 10). No grammar, dictionary or text collection was found.

\section{Primary Educational Resources}
No educational materials in Yameo exist in the search results. With no known speakers, a heritage curriculum would have to rely on the sparse historical record and on the living related language, Yagua.

\section{Linguistic and Reference Materials}
\begin{itemize}
\item Wikipedia, Yameo language: \url{https://en.wikipedia.org/wiki/Yameo_language}
\item Numeral list for Yameo (Loreto, Peru) from a numerals database: \url{https://lingweb.eva.mpg.de/channumerals/Yameo.htm}
\end{itemize}

\section{Religious and Community Resources}
No community organisation or religious resources specific to Yameo turned up in the searches. The related Yagua people are described in Wikipedia: \url{https://en.wikipedia.com/wiki/Yagua}

\section{Digital and Audio Resources}
No audio, video or corpus resources for Yameo turned up in the searches.

\section{Curriculum Development Context}
Yameo has no known speakers, and the surviving documentation found here is minimal. A heritage or revitalization curriculum could not use Yameo as a language of instruction. Realistic options are culture and history units about the Yameo and the Loreto region, and linking to Yagua language resources. Any reconstruction would need specialist linguistic work using the primary sources the Wikipedia articles cite, notably Mason (1950), which this search did not locate directly.

\end{document}